\usepackage{iftex}
\ifPDFTeX
  \usepackage[utf8]{inputenc}
  \usepackage[T1,T2A]{fontenc}
\else
  \usepackage{fontspec}
\fi
\usepackage[russian]{babel}
\usepackage{color}
\usepackage{calc}
\usepackage{graphicx}
\usepackage{epstopdf}
\usepackage{hyperref}
\hypersetup{unicode,colorlinks}
\usetheme[progressbar=head,numbering=fraction,block=fill]{metropolis}
\usepackage{minted}
\ifPDFTeX
  \usepackage{dejavu}
\else
  \setsansfont{DejaVu Sans}
  \setmonofont{DejaVu Sans Mono}
\fi
%\usepackage{adjustbox}  % Позволяет сузить куски кода (или текст) ровно настолько, чтобы уместиться в слайд
\usepackage{csquotes}
\usepackage{upquote}

\usemintedstyle{solarized-light}
% Для слайдов с блоками кода используйте \begin{frame}[fragile].
\setminted[lean4]{autogobble=true,breaklines=false}
\newminted[lean]{lean4}{fontsize=\small}
\newminted[leansmall]{lean4}{fontsize=\footnotesize}
\newminted[leantiny]{lean4}{fontsize=\scriptsize}
\newmintinline[leaninline]{lean4}{breaklines=true}

% Вывод #eval, сообщения Lean и состояния целей.
\newminted[leanoutput]{text}{autogobble=true,fontsize=\small}
\newminted[leanoutputsmall]{text}{autogobble=true,fontsize=\footnotesize}
\newminted[leanoutputtiny]{text}{autogobble=true,fontsize=\scriptsize}
\newmintinline[leanoutputinline]{text}{breaklines=true}

% Частые символы Lean в листингах при сборке pdfLaTeX.
\ifPDFTeX
  \usepackage{amssymb}
  \DeclareUnicodeCharacter{2192}{\ensuremath{\to}}
  \DeclareUnicodeCharacter{2190}{\ensuremath{\leftarrow}}
  \DeclareUnicodeCharacter{21D2}{\ensuremath{\Rightarrow}}
  \DeclareUnicodeCharacter{2194}{\ensuremath{\leftrightarrow}}
  \DeclareUnicodeCharacter{2200}{\ensuremath{\forall}}
  \DeclareUnicodeCharacter{2203}{\ensuremath{\exists}}
  \DeclareUnicodeCharacter{2227}{\ensuremath{\land}}
  \DeclareUnicodeCharacter{2228}{\ensuremath{\lor}}
  \DeclareUnicodeCharacter{00AC}{\ensuremath{\neg}}
  \DeclareUnicodeCharacter{2260}{\ensuremath{\ne}}
  \DeclareUnicodeCharacter{2264}{\ensuremath{\le}}
  \DeclareUnicodeCharacter{2265}{\ensuremath{\ge}}
  \DeclareUnicodeCharacter{2208}{\ensuremath{\in}}
  \DeclareUnicodeCharacter{22A2}{\ensuremath{\vdash}}
  \DeclareUnicodeCharacter{2115}{\ensuremath{\mathbb{N}}}
  \DeclareUnicodeCharacter{03B1}{\ensuremath{\alpha}}
  \DeclareUnicodeCharacter{03B2}{\ensuremath{\beta}}
  \DeclareUnicodeCharacter{27E8}{\ensuremath{\langle}}
  \DeclareUnicodeCharacter{27E9}{\ensuremath{\rangle}}
\fi

\vfuzz=20pt  % позволяет тексту дойти до номера слайда

\author{Алексей Романов}
\subtitle{Программирование и доказательства на Lean 4}
%\logo{}
\institute{МИЭТ}
\subject{Программирование и доказательства на Lean 4}
%\setbeamercovered{transparent}
%\setbeamertemplate{navigation symbols}{}
